% ============================================================
%  Part XIII — Strong and irregular verb tables
% ============================================================
\gpart{XIII}{Strong Verbs}

% ------------------------------------------------------------
\gsec{Strong and irregular verbs: beginnen -- fliegen}

\begin{gtbl}{colspec={X[1,l] X[0.9,l] X[0.9,l] X[1.2,l] X[1,l]}, rows={font=\footnotesize}, rowsep=1.8pt}
  Infinitiv & er/sie/es & Präteritum & Partizip II & Meaning \\
  beginnen  & beginnt   & begann  & hat begonnen  & to begin \\
  bieten    & bietet    & bot     & hat geboten   & to offer \\
  bitten    & bittet    & bat     & hat gebeten   & to ask for \\
  bleiben   & bleibt    & blieb   & ist geblieben & to stay \\
  brechen   & bricht    & brach   & hat gebrochen & to break \\
  bringen   & bringt    & brachte & hat gebracht  & to bring \\
  denken    & denkt     & dachte  & hat gedacht   & to think \\
  dürfen    & darf      & durfte  & hat gedurft   & to be allowed \\
  empfehlen & empfiehlt & empfahl & hat empfohlen & to recommend \\
  essen     & isst      & aß      & hat gegessen  & to eat \\
  fahren    & fährt     & fuhr    & ist gefahren  & to drive, go \\
  fallen    & fällt     & fiel    & ist gefallen  & to fall \\
  fangen    & fängt     & fing    & hat gefangen  & to catch \\
  finden    & findet    & fand    & hat gefunden  & to find \\
  fliegen   & fliegt    & flog    & ist geflogen  & to fly \\
\end{gtbl}

\begin{gnote}
The four columns you actually need to memorise are \textbf{Infinitiv,
er-Form, Präteritum, Partizip II}. Everything else in the paradigm is
predictable from them. The auxiliary is shown in the Partizip II column because
you can never guess it.
\end{gnote}

\begin{gnote}
Vowel patterns cluster. \textit{i--a--u} covers \textit{finden, singen,
trinken, springen}; \textit{ei--ie--ie} covers \textit{bleiben, scheinen,
schreiben, steigen}; \textit{ie--o--o} covers \textit{bieten, fliegen,
verlieren, ziehen}. Learn the pattern and a whole family comes with it.
\end{gnote}

% ------------------------------------------------------------
\gsec{Strong and irregular verbs: geben -- lügen}

\begin{gtbl}{colspec={X[1,l] X[0.9,l] X[0.9,l] X[1.2,l] X[1,l]}, rows={font=\footnotesize}, rowsep=1.8pt}
  Infinitiv & er/sie/es & Präteritum & Partizip II & Meaning \\
  geben     & gibt      & gab     & hat gegeben    & to give \\
  gefallen  & gefällt   & gefiel  & hat gefallen   & to please \\
  gehen     & geht      & ging    & ist gegangen   & to go \\
  gelingen  & gelingt   & gelang  & ist gelungen   & to succeed \\
  genießen  & genießt   & genoss  & hat genossen   & to enjoy \\
  geschehen & geschieht & geschah & ist geschehen  & to happen \\
  gewinnen  & gewinnt   & gewann  & hat gewonnen   & to win \\
  haben     & hat       & hatte   & hat gehabt     & to have \\
  halten    & hält      & hielt   & hat gehalten   & to hold \\
  hängen    & hängt     & hing    & hat gehangen   & to be hanging \\
  heißen    & heißt     & hieß    & hat geheißen   & to be called \\
  helfen    & hilft     & half    & hat geholfen   & to help \\
  kennen    & kennt     & kannte  & hat gekannt    & to know \\
  kommen    & kommt     & kam     & ist gekommen   & to come \\
  können    & kann      & konnte  & hat gekonnt    & to be able \\
  lassen    & lässt     & ließ    & hat gelassen   & to let \\
  laufen    & läuft     & lief    & ist gelaufen   & to run \\
  leihen    & leiht     & lieh    & hat geliehen   & to lend \\
  lesen     & liest     & las     & hat gelesen    & to read \\
  liegen    & liegt     & lag     & hat gelegen    & to lie \\
  lügen     & lügt      & log     & hat gelogen    & to tell a lie \\
\end{gtbl}

\begin{gnote}
\textit{kennen} and \textit{wissen} both mean \emph{to know}. Use
\textit{kennen} for people and places you are acquainted with, \textit{wissen}
for facts and for anything followed by a clause.
\end{gnote}

% ------------------------------------------------------------
\gsec{Strong and irregular verbs: messen -- sprechen}

\begin{gtbl}{colspec={X[1,l] X[0.9,l] X[0.9,l] X[1.2,l] X[1,l]}, rows={font=\footnotesize}, rowsep=1.8pt}
  Infinitiv & er/sie/es & Präteritum & Partizip II & Meaning \\
  messen    & misst     & maß     & hat gemessen    & to measure \\
  mögen     & mag       & mochte  & hat gemocht     & to like \\
  müssen    & muss      & musste  & hat gemusst     & to have to \\
  nehmen    & nimmt     & nahm    & hat genommen    & to take \\
  nennen    & nennt     & nannte  & hat genannt     & to name \\
  raten     & rät       & riet    & hat geraten     & to advise \\
  rufen     & ruft      & rief    & hat gerufen     & to call \\
  scheinen  & scheint   & schien  & hat geschienen  & to shine, seem \\
  schlafen  & schläft   & schlief & hat geschlafen  & to sleep \\
  schlagen  & schlägt   & schlug  & hat geschlagen  & to hit \\
  schließen & schließt  & schloss & hat geschlossen & to close \\
  schneiden & schneidet & schnitt & hat geschnitten & to cut \\
  schreiben & schreibt  & schrieb & hat geschrieben & to write \\
  schwimmen & schwimmt  & schwamm & ist geschwommen & to swim \\
  sehen     & sieht     & sah     & hat gesehen     & to see \\
  sein      & ist       & war     & ist gewesen     & to be \\
  singen    & singt     & sang    & hat gesungen    & to sing \\
  sitzen    & sitzt     & saß     & hat gesessen    & to sit \\
  sollen    & soll      & sollte  & hat gesollt     & to be supposed to \\
  sprechen  & spricht   & sprach  & hat gesprochen  & to speak \\
\end{gtbl}

\begin{gnote}
Modal verbs form their Partizip II regularly (\textit{gekonnt},
\textit{gemusst}) but you rarely see it. With a second verb they use the
\textbf{double infinitive} instead: \textit{Ich habe gehen \textbf{müssen}}, not
\textit{gemusst}.
\end{gnote}

% ------------------------------------------------------------
\gsec{Strong and irregular verbs: springen -- ziehen}

\begin{gtbl}{colspec={X[1,l] X[0.9,l] X[0.9,l] X[1.2,l] X[1,l]}, rows={font=\footnotesize}, rowsep=1.8pt}
  Infinitiv & er/sie/es & Präteritum & Partizip II & Meaning \\
  springen  & springt   & sprang  & ist gesprungen & to jump \\
  stehen    & steht     & stand   & hat gestanden  & to stand \\
  stehlen   & stiehlt   & stahl   & hat gestohlen  & to steal \\
  steigen   & steigt    & stieg   & ist gestiegen  & to climb \\
  sterben   & stirbt    & starb   & ist gestorben  & to die \\
  tragen    & trägt     & trug    & hat getragen   & to carry, wear \\
  treffen   & trifft    & traf    & hat getroffen  & to meet \\
  trinken   & trinkt    & trank   & hat getrunken  & to drink \\
  tun       & tut       & tat     & hat getan      & to do \\
  vergessen & vergisst  & vergaß  & hat vergessen  & to forget \\
  verlieren & verliert  & verlor  & hat verloren   & to lose \\
  wachsen   & wächst    & wuchs   & ist gewachsen  & to grow \\
  waschen   & wäscht    & wusch   & hat gewaschen  & to wash \\
  werden    & wird      & wurde   & ist geworden   & to become \\
  werfen    & wirft     & warf    & hat geworfen   & to throw \\
  wissen    & weiß      & wusste  & hat gewusst    & to know \\
  wollen    & will      & wollte  & hat gewollt    & to want \\
  ziehen    & zieht     & zog     & hat gezogen    & to pull \\
\end{gtbl}

\begin{gnote}
Verbs of motion and change take \textit{sein}: \textit{springen, steigen,
sterben, wachsen, werden}. Verbs that can take a direct object take
\textit{haben}, even when they feel like movement — \textit{Ich habe den Koffer
getragen}.
\end{gnote}

\begin{gnote}
Prefixed forms inherit everything from the base verb: \textit{verstehen} follows
\textit{stehen} (\textit{verstand}, \textit{verstanden}), \textit{aufstehen}
follows it too but takes \textit{sein}. Learn the 74 base verbs on these four
pages and several hundred compounds come free.
\end{gnote}
